\chapter{Prefácio}

Este livro tem uma ambição simples: uma única obra de química coerente,
que cubra tudo do primeiro ano da escola ao fim da graduação, escrita em
português para leitores de qualquer lugar do mundo.

O estilo é conciso e rigoroso: cursos construídos a partir de
definições, exemplos, proposições, métodos e esquemas, com um raciocínio
cuidadoso sempre que for acessível ao nível considerado, seguidos de
exercícios cuidadosamente escolhidos. Os resultados enunciados sem
dedução completa são marcados explicitamente como admitidos. As soluções
completas de todos os exercícios estão reunidas no final do livro. As
equações químicas estão sempre balanceadas, e todo número citado para
uma substância real vem de uma tabela de referência.

O curso é composto de quatro livros. O primeiro cobre os doze anos
escolares num único volume e nunca ensina a mesma noção duas vezes: cada
ideia é ensinada uma só vez, no ano em que ela se encaixa pela primeira
vez, e um capítulo posterior que precise dela começa com um breve
lembrete, \emph{O que você já sabe}, antes de ir mais longe. Os outros
três livros são os três anos da universidade. No primeiro livro, cada
ano escolar é uma parte, dividida em capítulos; cada livro universitário
é um ano, dividido em capítulos. Quanto mais cedo o ano, mais jovem é o
leitor a que ele se destina: mais ilustrações, mais passos detalhados e
exercícios mais suaves.

\section*{Como ler este livro}

Os enunciados são numerados em cada capítulo e codificados por cor: as
definições em azul, os teoremas em vermelho, as proposições e seus
parentes em laranja, e os métodos --- receitas curtas para as tarefas
habituais --- em verde. Os exercícios são graduados de ($\star$) para as
aplicações diretas a ($\star\star\star$) para os problemas mais
difíceis. Os quadros emoldurados contêm o que não faz parte do
raciocínio: um lembrete do que um capítulo anterior ensinou, a maneira
como um procedimento é realizado no laboratório, uma página de história,
ou os perigos de um produto.

\section*{Segurança}

A química se aprende com substâncias reais, e algumas delas são
perigosas. Os experimentos descritos neste livro são realizados num
laboratório escolar ou universitário, por um professor ou com ele, com
os equipamentos de proteção que o laboratório fornece. Nenhum deles é
feito para ser realizado em casa.

\section*{Contribuir}

O código-fonte deste livro está num repositório git público:
\begin{center}
\url{https://github.com/bvirrion/one-chemistry-book}
\end{center}
As contribuições --- novos capítulos, correções, melhores explicações,
exercícios adicionais --- são bem-vindas; veja \texttt{CONTRIBUTING.md}
na raiz do repositório.
